%! Equivalent Representations of Polynomial and Rational Expressions — Unit 1, Lesson 1.11 — Exit Ticket
\documentclass[10pt]{article}
\usepackage{precalculus-article}
\usepackage{precalculus-boxes}
\newcommand{\TallMath}[1]{$\displaystyle #1\rule[-1.4em]{0pt}{3.2em}$}

\begin{document}

\pageheader{Unit 1, Lesson 1.11}{Exit Ticket}
\namedateperiod

\vspace{0.14in}

\begin{enumerate}[itemsep=16pt]

\item A student needs to find the \textbf{vertical asymptotes and holes} of a rational function.
      Circle the form that is most useful, and explain why in one phrase.

      \medskip
      \textbf{A.} Factored form \qquad\qquad \textbf{B.} Standard form

      \vspace{0.06in}
      Because: \writeline

\item Perform the first step of polynomial long division to divide $x^2 + 3x + 5$ by $x + 2$.

      What is the first term of the quotient? \blank{3cm}

      After multiplying and subtracting, what is the new dividend? \blank{5cm}

\item Use Pascal's Triangle to expand $(x + 2)^4$.
      (Row 4 of Pascal's Triangle: $1 \quad 4 \quad 6 \quad 4 \quad 1$)

      \vspace{0.06in}
      $(x+2)^4 = $ \blank{9cm}

\end{enumerate}

\end{document}
